
faisceau tangent relatif & II App. II & 1.2.1\\
\((b)\)-famille & XIII & 1.12\\
famille ample & II & 2.2.4\\
famille de classes de faisceaux cohérents & XIII & 1.12\\
famille de classes de faisceaux cohérents limitée & XIII & 1.12\\
famille de classes de faisceaux cohérents limitée par un faisceau cohérent & XIII & 1.12\\
fibré cotangent relatif & II App. II & 1.2.1\\
fibré tangent relatif & II App. II & 1.2.1\\
fidèlement plat (morphisme de topos annelés) & I & 2.16.2\\
filtration de \(\lambda\)-algèbre augmentée & V & 3.10\\
\(\underline S\)-foncteur additif & I & 1.1.3\\
\(\underline S\)-foncteur exact & I & 1.1.3\\
formule de projection & IV & 2.11\\
groupe de Grothendieck & IV & 2.2\\
\(\lambda\)-homomorphisme & V & 2.1\\
homomorphisme de spécialisation & X App. & 7.3\\
\qquad''\hspace{8em}'' & X App. & 7.16\\
\(\lambda\)-idéal & V & 4.5\\
indice analytique & II App. II & 3.2.2\\
d'intersection complète (morphisme) & VIII & 1.1\\
\(n\)-isomorphisme & I & 1.4.3\\
libre homogène (module gradué) & VI & 1.2.2\\
localement stable par noyau d'épimorphisme & I & 1.2\\
localement quasi-relevable & I & 1.2\\
localement quasi-stable par noyau d'épimorphisme & I & 1.2\\
localement relevable & I & 1.2\\
nombre de Chern & X & 4.2\\
nombre d'intersection & X & 4.1\\
nombre d'intersection (pour une famille de faisceaux inversibles) & X & 4.4\\
numériquement équivalent à zéro (faisceau inversible) & XIII & 4.4\\
opérateur différentiel universel d'ordre \(n\) & II App. II & 1.2.1\\
